% chapters/07/section-02.tex
\section{Deteksi Tabrakan dengan Persegi Pembatas}
\label{sec:deteksi-tabrakan}

Bagian~2.3 telah menjelaskan gagasan persegi pembatas. Bagian ini memakai
gagasan itu melalui perkakas yang telah disediakan Pygame.

\subsection{Memeriksa Dua Persegi}

Pemeriksaan paling dasar adalah memeriksa dua persegi pembatas. Perintahnya
tersedia pada setiap \code{Rect}, sehingga tidak perlu dihitung sendiri.

\begin{kodekeluaran}
if player.rect.colliderect(platform.rect):
    # kedua persegi bersinggungan
\end{kodekeluaran}

Perintah \code{colliderect} mengembalikan nilai benar apabila kedua persegi
beririsan, dan salah apabila tidak. Dua persegi yang hanya
bersentuhan pada tepinya---misalnya sisi kanan yang satu tepat menyentuh sisi
kiri yang lain---dinilai \textbf{tidak} beririsan. Sifat itu berguna, sebab
benda yang sudah ditempelkan pada tepi bongkahan tidak akan dianggap
bertabrakan lagi pada bingkai berikutnya.

\subsection{Memeriksa Benda terhadap Wadah}

Ketika yang diperiksa adalah satu benda terhadap sekumpulan benda, Pygame
menyediakan perintah yang langsung menelusuri seluruh anggota wadah.

\begin{kodekeluaran}
hit_items = pygame.sprite.spritecollide(player, items, True)
\end{kodekeluaran}

Perintah itu mengembalikan daftar benda yang bersentuhan dengan pemain.
Argumen ketiga menentukan apakah benda yang bersentuhan langsung dikeluarkan
dari wadahnya. Nilai benar dipakai untuk kepingan yang diambil, sebab kepingan
yang telah diambil tidak diperlukan lagi. Nilai salah dipakai untuk rintangan,
sebab rintangan tetap ada di tempatnya setelah tersentuh.

Perhatikan nama variabel \code{hit_items} yang berbentuk daftar. Karena
hasilnya daftar, jumlah kepingan yang terambil dalam satu bingkai dapat
diketahui dengan \code{len()}. Cara itu lebih tepat daripada menambahkan satu
pada setiap pemeriksaan, sebab dalam satu bingkai pemain dapat menyentuh lebih
dari satu kepingan.

\subsection{Memeriksa Dua Wadah}

Apabila yang diperiksa adalah dua kumpulan benda, tersedia perintah ketiga.

\begin{kodekeluaran}
bertabrakan = pygame.sprite.groupcollide(bullets, enemies, True, True)
\end{kodekeluaran}

Perintah itu memeriksa setiap anggota wadah pertama terhadap setiap anggota
wadah kedua, dan mengembalikan kamus yang memetakan benda dari wadah pertama
ke daftar benda dari wadah kedua yang bersentuhan dengannya. Kedua argumen
penghapusan berlaku terpisah untuk masing-masing wadah.

Tabel~\ref{tab:perintah-tabrakan} merangkum ketiga perintah itu beserta
keperluannya.

\begin{table}[htbp]
\centering
\caption{Tiga perintah pemeriksaan tabrakan dan keperluannya}
\label{tab:perintah-tabrakan}
{\setlength{\tabcolsep}{4pt}\footnotesize%
\begin{tabularx}{\textwidth}{@{}L{5.0cm} L{3.2cm} Y@{}}
\toprule
\headrow \textbf{Perintah} & \textbf{Yang diperiksa} & \textbf{Dipakai ketika} \\
\midrule
\code{rect.colliderect(other)} & Satu persegi terhadap satu persegi &
Pemeriksaan diselesaikan sendiri, misalnya ketika menyelesaikan tabrakan sumbu
demi sumbu pada Bagian~\ref{sec:per-sumbu}. \\
\addlinespace
\code{spritecollide(benda, wadah, dokill)} &
Satu benda terhadap sekumpulan benda &
Pemain terhadap kepingan, rintangan, atau musuh. Perintah yang paling sering
dipakai. \\
\addlinespace
\code{groupcollide(wadah1, wadah2, kill1, kill2)} &
Sekumpulan benda terhadap sekumpulan benda &
Dua kumpulan yang keduanya banyak, misalnya peluru terhadap musuh. Hasilnya
kamus, bukan daftar. \\
\bottomrule
\end{tabularx}}
\end{table}

\subsection{Antara Mendeteksi dan Menyelesaikan}

Satu hal yang perlu ditegaskan sebelum melangkah lebih jauh: \textbf{mendeteksi
tabrakan dan menyelesaikannya adalah dua pekerjaan yang berbeda}. Seluruh
perintah pada bagian ini hanya menjawab pertanyaan ``apakah kedua benda
bersentuhan?''. Pertanyaan berikutnya---``dari sisi mana tabrakan terjadi, dan
apa yang harus dilakukan?''---belum terjawab.

Jawaban atas pertanyaan kedua itulah inti bab ini. Kedua
pertanyaan itu berbeda sifatnya: yang pertama dapat dijawab oleh pustaka,
sedangkan yang kedua bergantung pada aturan permainan Anda. Pada permainan
yang satu, menyentuh kepingan berarti skor bertambah; pada permainan yang lain,
menyentuh benda yang sama berarti kehilangan nyawa. Pustaka tidak dapat
mengetahuinya.
